% ============================================================================
% TECHNICAL APPENDIX 12A: Measuring HITL System Performance
% Formal definitions, calibration metrics, statistical inference,
% and experimental design for human-in-the-loop evaluation
% ============================================================================
% Source: /markdown/ch12/Ch12A_final.md
% ============================================================================

\chapter*{Technical Appendix 12A: Measuring HITL System Performance}
\addcontentsline{toc}{chapter}{Technical Appendix 12A: Measuring HITL Performance}

\begin{chapterquote}{Formal foundations for Chapter 12}
Measurement theory applied to human-in-the-loop systems: calibration metrics, pathway decomposition, inter-rater reliability, and the statistics of controlled experiments.
\end{chapterquote}

% ============================================================================
\section{Classification Metrics: Formal Definitions}
% ============================================================================

\subsection{The Confusion Matrix}

For a binary classifier, outcomes fall into four cells:

\begin{table}[htbp]
\centering
\caption{Binary classifier confusion matrix}
\label{tab:confusion-matrix}
\begin{tabular}{@{}lcc@{}}
\toprule
 & \textbf{Predicted Positive} & \textbf{Predicted Negative} \\
\midrule
\textbf{Actual Positive} & True Positive (TP) & False Negative (FN) \\
\textbf{Actual Negative} & False Positive (FP) & True Negative (TN) \\
\bottomrule
\end{tabular}
\end{table}

\subsection{Core Metrics}

\paragraph{Accuracy.}
\[
\text{Accuracy} = \frac{TP + TN}{TP + TN + FP + FN}
\]
Misleading under class imbalance.

\paragraph{Precision (Positive Predictive Value).}
\[
\text{Precision} = \frac{TP}{TP + FP}
\]
Of all cases flagged as positive, how many were correct? A precision-focused system minimizes false alarms.

\paragraph{Recall (Sensitivity, True Positive Rate).}
\[
\text{Recall} = \frac{TP}{TP + FN}
\]
Of all actually positive cases, how many were caught? A recall-focused system minimizes missed detections.

\paragraph{F1 Score.}
\[
F_1 = 2 \cdot \frac{\text{Precision} \times \text{Recall}}{\text{Precision} + \text{Recall}} = \frac{2 \cdot TP}{2 \cdot TP + FP + FN}
\]
The harmonic mean of precision and recall. The harmonic mean penalizes extreme imbalances between precision and recall more severely than the arithmetic mean.

\paragraph{$F_\beta$ Score.}
\[
F_\beta = (1 + \beta^2) \cdot \frac{\text{Precision} \times \text{Recall}}{(\beta^2 \cdot \text{Precision}) + \text{Recall}}
\]
When $\beta > 1$, recall is weighted more heavily (missing a positive is more costly). When $\beta < 1$, precision is weighted more heavily.

\subsection{HITL-Specific Pathway Decomposition}

Let $\mathcal{A}$ denote the set of auto-processed cases and $\mathcal{H}$ the set of human-reviewed cases, with $|\mathcal{A}| + |\mathcal{H}| = N$.

\paragraph{Pathway Error Rates.}
\[
\text{Error}_{\mathcal{A}} = \frac{FP_{\mathcal{A}} + FN_{\mathcal{A}}}{|\mathcal{A}|}, \qquad
\text{Error}_{\mathcal{H}} = \frac{FP_{\mathcal{H}} + FN_{\mathcal{H}}}{|\mathcal{H}|}
\]

\paragraph{System-Level Error Rate.}
\[
\text{Error}_{\text{system}} = \frac{|\mathcal{A}|}{N} \cdot \text{Error}_{\mathcal{A}} + \frac{|\mathcal{H}|}{N} \cdot \text{Error}_{\mathcal{H}}
\]

A system can show improving aggregate error while $\text{Error}_{\mathcal{A}}$ increases, if the routing threshold is raised (more cases auto-processed). Tracking pathway-specific error rates separately detects this divergence.

\paragraph{Human Value-Added.}
\[
\Delta_H = \text{Error}_{\mathcal{H}}^{(\text{auto-only})} - \text{Error}_{\mathcal{H}}^{(\text{human-reviewed})}
\]
Estimated counterfactual: how much higher would the error rate on human-reviewed cases have been if the model had decided them autonomously?

% ============================================================================
\section{Calibration Metrics}
% ============================================================================

\subsection{Why Calibration Matters for HITL}

A HITL system routes cases to humans based on model confidence. If confidence is miscalibrated---if an 80\% confidence score corresponds to 60\% accuracy---the routing threshold is miscalibrated too.

\paragraph{Definition: Calibration.} A probabilistic classifier $f$ is perfectly calibrated if, for all $p \in [0,1]$:
\[
P\bigl(\hat{y} = y \mid f(x) = p\bigr) = p
\]

\subsection{Expected Calibration Error (ECE)}

Partition predictions into $M$ equally-spaced confidence bins. Let $B_m$ denote the set of cases in bin $m$, with $|B_m|$ elements:
\[
\text{acc}(B_m) = \frac{1}{|B_m|} \sum_{i \in B_m} \mathbf{1}[\hat{y}_i = y_i], \qquad
\text{conf}(B_m) = \frac{1}{|B_m|} \sum_{i \in B_m} \hat{p}_i
\]

\[
\text{ECE} = \sum_{m=1}^{M} \frac{|B_m|}{n} \left| \text{acc}(B_m) - \text{conf}(B_m) \right|
\]

A perfectly calibrated model has $\text{ECE} = 0$. Values below 0.02 are considered good; above 0.05 warrant intervention \autocite{guo2017calibration}.

\subsection{Maximum Calibration Error (MCE)}

ECE is an average; a model can have low ECE while severely miscalibrated in high-confidence bins where routing decisions are made:
\[
\text{MCE} = \max_m \left| \text{acc}(B_m) - \text{conf}(B_m) \right|
\]

For HITL systems, MCE in the high-confidence bins should be tracked separately from the overall ECE.

\subsection{Calibration Methods}

\paragraph{Temperature Scaling \autocite{guo2017calibration}.}
A single learned parameter $T > 0$ scales the logits before the softmax:
\[
\hat{p}_i = \text{softmax}(z_i / T)
\]
When $T > 1$, confidence is reduced. $T$ is optimized on a held-out validation set by minimizing negative log-likelihood.

\paragraph{Platt Scaling.}
Fits a logistic regression $P(y=1 \mid f) = \sigma(af + b)$ on a validation set, where $(a, b)$ are learned parameters.

\paragraph{Isotonic Regression.}
Non-parametric monotonic calibration. More flexible than temperature scaling but requires a larger validation set to avoid overfitting. Preferred when calibration errors are concentrated in specific confidence regions.

% ============================================================================
\section{Confidence Intervals for HITL Metrics}
% ============================================================================

\subsection{Wilson Score Interval for Proportions}

For an observed proportion $\hat{p} = k/n$, the Wilson score interval:
\[
\left[ \frac{\hat{p} + z^2/(2n) \pm z\sqrt{\hat{p}(1-\hat{p})/n + z^2/(4n^2)}}{1 + z^2/n} \right]
\]
where $z = 1.96$ for a 95\% interval. Preferred over the Wald interval when $\hat{p}$ is near 0 or 1.

\subsection{Bootstrap Confidence Intervals}

F1, ECE, and inter-rater agreement are nonlinear functions of the data. Closed-form intervals are unreliable. The bootstrap:
\begin{enumerate}
    \item Draw $B = 10{,}000$ bootstrap samples (with replacement) from $n$ observations.
    \item Compute the metric on each: $\theta_1^*, \ldots, \theta_B^*$.
    \item The 95\% interval is the 2.5th and 97.5th percentiles of $\{\theta_b^*\}$.
\end{enumerate}

\subsection{Minimum Detectable Effect}

Required sample size to detect error-rate difference $\delta$ with power $1-\beta$ at significance $\alpha$:
\[
n = \frac{(z_{\alpha/2} + z_\beta)^2 \cdot 2\bar{p}(1-\bar{p})}{\delta^2}
\]
where $\bar{p}$ is the pooled error rate. For a 5\% baseline and 1 percentage point improvement, 80\% power requires $n \approx 2{,}600$ cases per arm.

% ============================================================================
\section{Inter-Rater Reliability}
% ============================================================================

\subsection{Cohen's Kappa}

Raw agreement is inflated by chance. Kappa corrects for this:
\[
\kappa = \frac{p_o - p_e}{1 - p_e}
\]
where $p_o$ is observed agreement and $p_e = \sum_k (n_{k\cdot}/n)(n_{\cdot k}/n)$ is expected chance agreement.

Benchmarks (Landis \& Koch, 1977): $\kappa < 0.20$ slight; $0.21$--$0.40$ fair; $0.41$--$0.60$ moderate; $0.61$--$0.80$ substantial; ${>}0.80$ nearly perfect.

\subsection{Fleiss' Kappa for Multiple Reviewers}

For $k > 2$ reviewers, Fleiss' generalization:
\[
\kappa_F = \frac{\bar{P} - \bar{P}_e}{1 - \bar{P}_e}
\]

\subsection{Intraclass Correlation Coefficient (ICC)}

For continuous annotations (e.g., severity scores), two-way mixed ICC:
\[
\text{ICC}(2,1) = \frac{MS_R - MS_E}{MS_R + (k-1)MS_E + k(MS_C - MS_E)/n}
\]
where $MS_R$, $MS_C$, $MS_E$ are mean squares from a two-way ANOVA decomposition.

% ============================================================================
\section{Statistical Significance in A/B Tests}
% ============================================================================

\subsection{Two-Proportion Z-Test}

Null hypothesis $H_0: p_A = p_B$. Test statistic:
\[
z = \frac{\hat{p}_A - \hat{p}_B}{\sqrt{\hat{p}(1-\hat{p})(1/n_A + 1/n_B)}}
\]
where $\hat{p} = (k_A + k_B)/(n_A + n_B)$. Reject $H_0$ at $\alpha = 0.05$ when $|z| > 1.96$.

\subsection{Complications in HITL Experiments}

\begin{keyconcept}[Three HITL-Specific Complications]
\begin{itemize}
    \item \textbf{Non-independence:} Cases reviewed by the same reviewer are correlated. Use cluster-robust variance estimation with reviewers as clusters.
    \item \textbf{SUTVA violations:} Reviewer behavior can change based on queue length, which is affected by arm assignment. Minimize by balancing queue lengths across arms.
    \item \textbf{Learning effects:} If the model updates during the experiment, freeze model weights for the experimental period or use pre-registered stopping rules.
\end{itemize}
\end{keyconcept}

\subsection{CUPED: Variance Reduction Using Pre-Experiment Data}

CUPED reduces variance by regressing out baseline covariates \autocite{deng2013improving}:
\[
Y_{\text{CUPED}} = Y - \theta(X - \mathbb{E}[X])
\]
where $\theta = \text{Cov}(Y, X)/\text{Var}(X)$ is estimated from pre-experiment data. Can reduce variance by 30--50\%, roughly halving required sample size.

% ============================================================================
\section{Feedback Loop Latency Measurement}
% ============================================================================

\subsection{Defining Latency}

Let $t_{\text{decision}}$ be the time a human reviewer finalizes a decision, and $t_{\text{incorporated}}$ be the time that decision is reflected in the model's predictions:
\[
L_i = t_{\text{incorporated},i} - t_{\text{decision},i}, \qquad
\bar{L} = \frac{1}{|\mathcal{H}|} \sum_{i \in \mathcal{H}} L_i
\]

\subsection{Measuring Feedback Value}

The value of the feedback loop is the improvement attributable to human corrections:
\[
V_{\text{feedback}} = \text{Error}_t^{(\text{without feedback})} - \text{Error}_{t+\Delta}^{(\text{with feedback})}
\]
Estimating the counterfactual requires freezing model snapshots at multiple time points and evaluating each against new ground truth---operationally complex but the only rigorous approach.

% ============================================================================
\section{Exercises}
% ============================================================================

\begin{Exercise}[Calibration Assessment]
A model has predictions in 10 confidence bins with accuracy values: \{0.05, 0.11, 0.18, 0.30, 0.42, 0.55, 0.63, 0.74, 0.82, 0.90\} (equal numbers per bin, $n=1{,}000$ total).
\begin{enumerate}[(a)]
    \item Compute the ECE.
    \item Describe the reliability diagram shape.
    \item Assess calibration quality; is temperature scaling appropriate?
    \item What does this calibration imply for HITL routing at a threshold of 0.80?
\end{enumerate}
\end{Exercise}

\begin{Exercise}[Sample Size Determination]
A HITL system has a current error rate of 8\%. A proposed redesign is expected to reduce this to 6\%.
\begin{enumerate}[(a)]
    \item Compute the required sample size per arm at 80\% power, $\alpha = 0.05$.
    \item Recompute for 90\% power.
    \item At a throughput of 500 human-reviewed cases per day, how long would each experiment take? What practical challenges arise at this timescale?
\end{enumerate}
\end{Exercise}

\begin{Exercise}[Goodhart's Law Simulation]
A model routes the bottom 20\% of confidence cases to human review. Original reviewer accuracy: 92\%. Gaming increases throughput 3$\times$ but reduces accuracy to 65\%. System auto-processes the remaining 80\% at a 3\% error rate.
\begin{enumerate}[(a)]
    \item Compute the system-level error rate before and after gaming using the pathway decomposition from \cref{sec:divergence-trap}.
    \item Explain why aggregate accuracy might appear stable to management while the system degrades.
    \item What single monitoring metric would most reliably detect gaming within two weeks?
\end{enumerate}
\end{Exercise}

\begin{Exercise}[Inter-Rater Reliability]
Two reviewers independently label 100 borderline cases. Reviewer A: 60 positive, 40 negative. Reviewer B: 55 positive, 45 negative. Agreement: 50 both-positive, 25 both-negative.
\begin{enumerate}[(a)]
    \item Compute Cohen's kappa.
    \item Interpret the level of agreement.
    \item What trend in kappa over time would constitute a warning signal? What might cause it?
\end{enumerate}
\end{Exercise}

\begin{Exercise}[Feedback Loop Quality Control]
Design a quality control procedure for incoming reviewer decisions before model training. The system processes 10,000 cases per day, routes 15\% to human review, and targets 24-hour feedback loop latency.
\begin{enumerate}[(a)]
    \item Specify four statistics to compute on each incoming daily batch.
    \item Define alert thresholds for each statistic.
    \item Describe the decision rule: when to proceed, when to hold for audit, when to reject.
    \item What is the cost of a false hold (rejecting a good batch)? Of a false pass (accepting a bad batch)?
\end{enumerate}
\end{Exercise}

\bigskip
\begin{center}
\rule{0.5\textwidth}{0.4pt}
\end{center}
\bigskip

\noindent\textit{This appendix provides the mathematical foundations for Chapter~12's measurement framework. The exercises are calibrated so that Exercises~12A.1--12A.4 are accessible to readers with introductory statistics, while Exercise~12A.5 requires systems-level thinking appropriate for practitioners and advanced students.}
